\section{Using the full first-index modulus}
\label{sec:indexmodulus}

The signed Pell family admits a separate decrease in degree. The complete
circuits supplied with this section have the following exact sizes; every
unknown ranges over all integers.
\begin{center}
\begin{tabular}{rrrr}
Unknowns&Old degree/operations&New degree&New operations\\\hline
14&86/112&82&113\\
14&94/110&90&111\\
13&156/125&152&126\\
11&300/148&298&150
\end{tabular}
\end{center}
These are degree--cost tradeoffs. This separate branch retains the repository's native/K215 specialization; its fixed coefficient tuple must not be confused with the direct quartic, which has small literal coefficients. Their six fixed native coefficients are
the complete source34 tuple supplied in the component, rather than new
unknowns or an unnamed universal polynomial.

\subsection{The change in the arithmetic kernel}
The inherited signed front defines
\begin{align*}
 J&=6j+1,&q&=\beta J+1,&W&=12w+4,\\
 X&=qW,&Y&=v,&&\\
 F&=(K_0+W)C-z(q-1),&&&&\\
 R&=[q(q-F)-C](q^2-1)+(M_C+qM_F)J.&&&&
\end{align*}
Write
\[
 E=XY,\quad a=E+Y,\quad A=a+2,\quad H=4a+3,
 \quad \Delta=A^2-1.
\]
Replace the old index coordinate and magnitude budget by
\begin{align}
 \ell&=XY^2,&k&=R+1+h\ell,&\eta&=c-Yk,\label{eq:rebasedk}\\
 D_e&=(\mathrm{cone}_A F-q)^2+(\mathrm{cone}_B C-q)^2,&&&\notag\\
 B_0&=k^2(q^2-D_e)-4q^2\eta^2,&&&\notag\\
 \mathcal B_6&=Y^2B_0-k^2q^6.&&&\label{eq:rebasebudget}
\end{align}
The predecessors used $k=R+1+hXY$ and
$\mathcal B_8=Y^2B_0-k^2q^8$. The first Pell radicand is unchanged
as a function of the semantic values:
\[
 D_t=1+\ell(\ell+1)k^2.
\]

At any new zero the guard gives $\mathcal B_6>0$. Thus $kY\ne0$,
$B_0>0$, and, since $B_0\le k^2q^2$,
\begin{equation}\label{eq:rebasemagnitude}
 |Y|>|q|^2.
\end{equation}
The centered-disk argument recovers the common signs of $q,F,C,q-F-C$
and $|c/k-Y|<1/2$. Put $q_0=|q|$, $x=|X|$, $m=|Y|$.
The authentic coefficient bounds supply
\[
 q_0\ge16,\qquad x\ge4q_0,\qquad
 q_0(q_0^2-1)<R<q_0^4-2.
\]
The larger modulus now satisfies
\begin{equation}\label{eq:rebasedmodulus}
 |\ell|=xm^2>4q_0^5>3R+1.
\end{equation}
This replaces the predecessor's early requirement $|XY|>3R+1$.

\subsection{Why the first index is still recovered exactly}
For the Pell ordinate sequence
$\psi_b(0)=0$, $\psi_b(1)=1$,
$\psi_b(n+2)=2b\psi_b(n+1)-\psi_b(n)$, induction gives
\[
 \psi_{2\ell+1}(n)\equiv n\pmod{2\ell}.
\]
The unchanged first/main Pell classification yields
\[
 |k|=2\psi_{|2\ell+1|}(n_0),\qquad
 |c|=\psi_{|A|}(p),\qquad p>n_0\ge1.
\]
The inequality $p>n_0$ follows from $|2\ell+1|>|A|$ and
$|c|>|k|$, both consequences of the centered strip and
\eqref{eq:rebasemagnitude}. The literal index definition gives
\[
 \epsilon\,2n_0\equiv R+1\pmod{|\ell|},\qquad
 \epsilon=\operatorname{sgn}(k)\operatorname{sgn}(X)^{n_0-1}.
\]
By \eqref{eq:rebasedmodulus}, the positive branch has
$n_0\ge(R+1)/2$; the negative branch has $n_0>R$.
These are exactly the lower index bounds needed by the retained signed
strong-rank and auxiliary step-down lemmas. They recover $p=R$.
The negative branch is then impossible, and the congruence cannot wrap:
\[
 2n_0=R+1.
\]

The signed ratio argument needs only $x\ge64$, $m\ge6$, this index
equality and the centered strip. It consequently gives, for $R=2t+1$,
\[
 m>\frac{(x-3)^t}{2}-1,\qquad |H|>x(x-3)^t.
\]
The inherited whole-product projection now forces
$X=2^R$ and the correct signs, and finally recovers
$Y=(2v_{\rm scale}+1)q^3$. The new soundness proof rejoins the earlier
$B_0$-only source, with
\[
 h_{\rm old}=Yh_{\rm new}.
\]
In particular it does \emph{not} infer $\mathcal B_8>0$ from
$\mathcal B_6>0$ at an arbitrary zero. Fresh ternary positivity witnesses
restore $B_0>0$ in the earlier source. This explicit distinction avoids a
circular use of the old stronger magnitude guard.

Conversely, canonical old witnesses have
$k=2\psi_{2\ell+1}((R+1)/2)$. Hence $4\ell$ divides $k-R-1$ and the
new $h=(k-R-1)/\ell$ is integral. Their margin gives
\[
 \mathcal B_6=\mathcal B_8+k^2(q^8-q^6)>\mathcal B_8>0.
\]
A fresh ternary representation completes the new witnesses. Thus both
directions of the calibrated zero-existence equivalence hold.

\subsection{A complete fourteen-integer polynomial of degree82}
The only unknowns are
\[
 j,C,z,h,r,v,w,T,c,n,f,p_2,p_1,s\in\mathbb Z.
\]
Use the preceding front and define the following polynomial abbreviations:
\begin{align*}
 L&=\ell k,& I&=(2r+1)c^2,&Q&=\Delta(f^2-1),\\
 V&=f^2(cT-R)-c,& M&=2q^3,&U&=MH,\\
 K_{\rm lo}&=X(2X-3c)-2,
 &Z&=Un+MK_{\rm lo}+H(Y-q^3),\\
 A_c&=(Z+Uk)^2+4U(LZ-U),
 &S_f&=f^2-\Delta I^2,\\
 B_g&=\mathcal B_6-p_1^2-p_2^2,
 &G&=4B_g-3,\\
 \delta&=Q-1,&C_y&=QV^2-1,\\
 K_s&=\delta\bigl(s(s+1)-B_g+1\bigr)-C_y,
 &N_s&=K_s^2-G\delta C_y.
\end{align*}
The equation is
\begin{equation}\label{eq:integer82}
 \boxed{(A_c+1)S_f(N_s+1)-1=0.}
\end{equation}
The first factor cannot be $-1$ because $A_c$ is a square modulo4.
The second cannot be $-1$ because $\Delta$ is0 or3 modulo4.
Their product equation therefore forces $A_c=0$, $S_f=1$, $N_s=0$.

The retained primitive odd-carrier lemma follows from
\begin{align*}
 &16\left([d(s(s+1)-b)-e]^2-(4b+1)de\right)\\
 &\qquad=[d((2s+1)^2+4b+1)-4e]^2
       -4d^2(2s+1)^2(4b+1).
\end{align*}
When $d\ne0$, its zero is equivalent to integral witnesses $g,u$ with
$g^2=4b+1$, $du^2=e$, and $2s+1=g+2u$.
Apply it to $b=B_g-1$, $d=\delta$, $e=C_y$.
The strong equation makes $Q=\Delta^2I^2$, so $\delta\ne0$.
Then $p_3=(g-1)/2$ and $y=u-V$ restore the omitted positivity and
auxiliary coordinates. The inverse assignment is $s=p_3+V+y$.
This proves the exact algebraic reduction to the changed kernel.

The complete actual degree is82. One certified nonzero coefficient is
\[
 [j^{47}w^9v^{10}z^5f^4r^2c^4n]\,P_{82}
       =-2^{75}3^{56}\beta^{47}.
\]
Every fixed coefficient has its full finite reconstruction in the supplied
component. The complete proof of the degree90 and degree152 companions
uses the same changed kernel and retained parity-root mergers.

\subsection{The eleven-integer degree298 endpoint}
The eleven-integer linear-cofactor family uses the same new
$k,\mathcal B_6$. Its incoming degree300 center can be improved by the
single integral correction
\[
 \Psi_{\rm new}=\Psi_{\rm old}-16rc^4U_s^2T,
 \qquad U_s=rc^2\bigl(1+(A^2-1)c^2\bigr).
\]
At this scalar family's final Q41/native handoff, explicitly set
$h_{\rm old}=Yh_{\rm new}$, so that the retained source computes
$R+1+h_{\rm old}XY=R+1+h_{\rm new}XY^2=k$; the other restored outer
coordinates and fresh strong witnesses follow the retained scalar
construction.
The complete circuit and full ordinary proof are in the signed component.
For the exact cancellation audit, assign weights $(3,1,1,1,5)$ to
$(A,c,r,T,R)$. The pre-auxiliary center remainder drops from degree82 to80,
and the full center remainder from110 to108. The norm's squared inner
term consequently has degree296. The weaker budget lowers the other
norm term from degree300 to298, which now dominates. The actual
coefficient
\[
 [j^{136}w^{46}v^{50}z^6s^6r^{12}c^{40}T^2]\,P_{298}
       =2^{272}3^{182}\beta^{136}
\]
is positive and nonzero. The independent sparse-ring checker verifies all
924 coefficients of the full center remainder, every circuit cut and the
attained degree at the actual coefficient tuple. These finite checks
authenticate the algebraic artifact; the unbounded Pell and logical
calibration remain the explicitly identified ordinary proof inputs.
